% © 2026 Ing. David Jaroš — CC BY-NC-ND 4.0
%
% This work is licensed under a Creative Commons Attribution-NonCommercial-NoDerivatives
% 4.0 International License (CC BY-NC-ND 4.0).
%
% License History: Earlier drafts (up to v0.3) were released under CC BY 4.0.
% From v0.4 onward, all material is released under CC BY-NC-ND 4.0 to protect
% the integrity of the theoretical work during ongoing academic development.

% UBT-AI-PROVENANCE-BEGIN
% schema: ubt-ai-provenance/v1
% tier: C_working
% ai_assistance: disclosed
% human_review: risk-based
% editorial_responsibility: Ing. David Jaroš
% policy: ../../AI_PROVENANCE.md
% notice: Working material; exhaustive human review is not claimed.
% UBT-AI-PROVENANCE-END

\documentclass[11pt,a4paper]{article}
\usepackage{amsmath,amssymb,amsthm}
\usepackage{booktabs}
\usepackage{geometry}
\usepackage{hyperref}
\geometry{margin=2.5cm}

\newtheorem{proposition}{Proposition}[section]
\newtheorem{remark}[proposition]{Remark}
\newtheorem{lemma}[proposition]{Lemma}

\title{Charge Quantization and Topology in UBT\\
\large Target 4 — fix\_or\_reject\_U1\_coupling\_normalization}
\author{Ing.\ David Jaroš}
\date{2026-05-10}

\begin{document}
\maketitle

\begin{abstract}
We analyze whether UBT topology forces electric charge quantization.  We examine
winding number quantization on $S^1_\psi$, holonomy conditions, allowed charges
from single-valuedness of $\Theta$, quark fractional charges, and the role of
the Weinberg angle.  Verdict: \textbf{CONDITIONAL} — integer charge
quantization follows from topology; fractional charges remain conditional.
\end{abstract}

% ----------------------------------------------------------------
\section{Winding number quantization on $S^1_\psi$}
% ----------------------------------------------------------------

The imaginary-time coordinate $\psi$ is compactified on a circle
$S^1_\psi$ of circumference $2\pi R_\psi$.  The fundamental field $\Theta$
must satisfy
\begin{equation}
  \Theta(x,\psi+2\pi R_\psi)=\Theta(x,\psi).
\end{equation}
Under a U(1) gauge transformation with parameter $\chi(\psi)$, the field
transforms as
\begin{equation}
  \Theta\;\longmapsto\; e^{i\chi(\psi)}\Theta.
\end{equation}
For $\Theta$ to be single-valued after one full circuit of $S^1_\psi$,
\begin{equation}
  \chi(2\pi R_\psi)-\chi(0) = 2\pi n,\qquad n\in\mathbb{Z}.
\end{equation}
This is the winding number condition.

% ----------------------------------------------------------------
\section{Holonomy and allowed charges}
% ----------------------------------------------------------------

The holonomy of the gauge connection around $S^1_\psi$ is
\begin{equation}
  \mathcal{W}=\exp\!\left(iq\oint_{S^1_\psi}A_\psi\,d\psi\right),
\end{equation}
where $q$ is the charge of the field under U(1).  Single-valuedness of $\Theta$
under parallel transport requires
\begin{equation}
  \mathcal{W} = 1,\qquad\text{i.e.,}\qquad
  q\oint_{S^1_\psi}A_\psi\,d\psi = 2\pi n,\qquad n\in\mathbb{Z}.
\end{equation}

\begin{proposition}[Charge quantization in integer units]
\label{prop:charge-quant}
If the minimal gauge flux through $S^1_\psi$ is normalized to the unit flux
quantum
\begin{equation}
  \Phi_0 := \oint_{S^1_\psi}A_\psi^{(1)}\,d\psi = 2\pi
  \quad\text{(unit flux quantum)},
\end{equation}
then the allowed charges of $\Theta$ are
\begin{equation}
  q = n \in \mathbb{Z}.
\end{equation}
\end{proposition}

\begin{remark}[CONDITIONAL status]
Proposition \ref{prop:charge-quant} requires that the unit flux quantum
$\Phi_0=2\pi$ is the minimal allowed flux.  This is the standard Dirac
quantization condition and holds if the gauge field is defined on a
simply-connected space.  On $S^1_\psi$, the flux is a Wilson line modulus and
is not a priori constrained to $2\pi$; the normalization $\Phi_0=2\pi$ is an
additional convention choice (equivalent to choosing $e_5=1$).
\end{remark}

% ----------------------------------------------------------------
\section{Derivation of allowed charges from single-valuedness}
% ----------------------------------------------------------------

For a field $\Theta$ with U(1) charge $q$ under the right-phase action
$\Theta\to e^{iq\chi}\Theta$, single-valuedness on $S^1_\psi$ with a
background holonomy $\mathcal{W}=e^{i\phi}$ requires
\begin{equation}
  e^{iq\phi}=1 \quad\Rightarrow\quad q\phi = 2\pi m,\quad m\in\mathbb{Z}.
\end{equation}
If only integer holonomies are allowed (which follows from Proposition
\ref{prop:charge-quant} above), then $\phi=2\pi$ and $q=m\in\mathbb{Z}$.

The fundamental biquaternion field $\Theta$ has charge $q_\Theta=1$ (derived in
Target 1).  Therefore:
\begin{itemize}
  \item $\Theta$: $q=1$ (electron-like, integer charge).
  \item $\Theta^n$: $q=n$ (integer multiples).
  \item Composites of $\Theta$ fields carry integer charges.
\end{itemize}

% ----------------------------------------------------------------
\section{Quark fractional charges and their origin}
% ----------------------------------------------------------------

Quarks carry charges $q=+2/3$ (up-type) and $q=-1/3$ (down-type), which are
not integers in units of the electron charge.

\subsection{Route via electroweak embedding}

The physical electric charge is $Q=T_3+Y/2$, where $T_3$ is the third component
of weak isospin and $Y$ is weak hypercharge.  The fractional charges of quarks
arise from their non-trivial hypercharge assignments:
\begin{equation}
  Y_q=\frac{1}{3}\quad\text{(quark doublets)},
  \qquad
  Y_u=\frac{4}{3},\quad Y_d=-\frac{2}{3}\quad\text{(quark singlets)}.
\end{equation}
These are rational numbers, not derived from the U(1)$_\psi$ topology alone.
They represent the embedding of quarks in the Standard Model gauge group
$\mathrm{SU}(3)_c\times\mathrm{SU}(2)_L\times\mathrm{U}(1)_Y$.

\begin{center}
\begin{tabular}{@{}lllll@{}}
\toprule
Field & $T_3$ & $Y/2$ & $Q=T_3+Y/2$ & Origin in UBT \\
\midrule
Electron $e^-$ & $-1/2$ & $-1/2$ & $-1$ & CONDITIONAL \\
Up quark $u$ & $+1/2$ & $+1/6$ & $+2/3$ & CONDITIONAL \\
Down quark $d$ & $-1/2$ & $+1/6$ & $-1/3$ & CONDITIONAL \\
\bottomrule
\end{tabular}
\end{center}

\subsection{Route via GUT embedding}

In a GUT-style embedding $\mathrm{SU}(5)\supset\mathrm{SU}(3)\times\mathrm{SU}(2)\times\mathrm{U}(1)$,
fractional charges $1/3$ arise from the GUT generator normalization.  This
requires a specific GUT group choice that is not currently derived from
$\mathcal{B}=\mathbb{C}\otimes\mathbb{H}$.

The biquaternion algebra $\mathcal{B}\cong\mathrm{Mat}(2,\mathbb{C})$ gives
$\mathrm{SU}(2)_L\times\mathrm{U}(1)_Y$ from the left factor automorphisms and
$\mathrm{SU}(3)_c$ from the involution structure.  The combination required for
$\mathrm{SU}(5)$ is not manifest.

\textbf{Status: CONDITIONAL (fractional charges require hypercharge assignments
imported from Standard Model or GUT).}

% ----------------------------------------------------------------
\section{Weinberg angle and U(1)$_Y$ to U(1)$_{\rm EM}$ mixing}
% ----------------------------------------------------------------

The physical photon is
\begin{equation}
  A_\mu=\sin\theta_W W^3_\mu+\cos\theta_W B_\mu,
\end{equation}
where $\theta_W$ is the Weinberg angle.  The electromagnetic coupling is
\begin{equation}
  e=g\sin\theta_W=g'\cos\theta_W.
\end{equation}
The Weinberg angle is related to the ratio of gauge couplings:
\begin{equation}
  \tan\theta_W=\frac{g'}{g}.
\end{equation}
This ratio is not fixed by the biquaternion algebra alone; it requires an
additional input (e.g., coupling unification at a GUT scale).

\textbf{Status: CONDITIONAL (Weinberg angle not derived from $\mathcal{B}$
alone).}

% ----------------------------------------------------------------
\section{Summary and verdict}
% ----------------------------------------------------------------

\begin{center}
\begin{tabular}{@{}lll@{}}
\toprule
Item & Status & Notes \\
\midrule
Single-valuedness on $S^1_\psi$ & DERIVED & Topological requirement \\
Integer charge quantization (unit flux) & CONDITIONAL & Requires $\Phi_0=2\pi$ convention \\
$q_\Theta=1$ from right-phase action & DERIVED & Algebraic, Target 1 \\
Integer charges for composites & DERIVED & From $q_\Theta=1$ \\
Quark charge $2/3$ & CONDITIONAL & Requires SM hypercharge assignment \\
Quark charge $-1/3$ & CONDITIONAL & Requires SM hypercharge assignment \\
Weinberg angle & CONDITIONAL & Ratio $g'/g$ not fixed by $\mathcal{B}$ \\
\bottomrule
\end{tabular}
\end{center}

\paragraph{Overall verdict.}
\textbf{CONDITIONAL.}
Integer charge quantization follows from the single-valuedness of $\Theta$ on
$S^1_\psi$, given the standard flux normalization convention.  The electron
carries unit charge by construction ($q_\Theta=1$).  Quark fractional charges
require Standard Model hypercharge assignments that are not derived from the
biquaternion algebra alone.  The Weinberg angle mixing remains an open parameter.

\begin{remark}[Potential path to fractional charges]
If a full $\mathrm{SU}(5)$ or $\mathrm{SO}(10)$ GUT structure were derived
from $\mathcal{B}$, fractional charges would follow automatically.  The
involution structure giving $\mathrm{SU}(3)_c$ is a promising starting point,
but the extension to full GUT has not been established.
\end{remark}

\end{document}
